% p411_1; source PDF page 411; manually verified editable table.
\begin{tikzpicture}[x=1cm,y=-1cm,inner sep=0pt,outer sep=0pt]
\path[use as bounding box] (0,0) rectangle (16.000,17.325);
\draw[line width=.4pt] (0.000,0.000) rectangle (3.100,0.865);
\node[align=center,text width=2.840cm,font={\linespread{1}\fontsize{10}{13.50}\selectfont\bfseries}] at (1.550,0.432) {误解的类型};
\draw[line width=.4pt] (3.100,0.000) rectangle (7.300,0.865);
\node[align=center,text width=3.940cm,font={\linespread{1}\fontsize{10}{13.50}\selectfont\bfseries}] at (5.200,0.432) {具体情形};
\draw[line width=.4pt] (7.300,0.000) rectangle (16.000,0.865);
\node[align=center,text width=8.440cm,font={\linespread{1}\fontsize{10}{13.50}\selectfont\bfseries}] at (11.650,0.432) {范例};
\draw[line width=.4pt] (0.000,0.865) rectangle (3.100,7.995);
\node[align=left,text width=2.840cm,font={\linespread{1}\fontsize{10}{13.50}\selectfont}] at (1.550,4.430) {误解一：只要客观上具备数额，就定罪量刑};
\draw[line width=.4pt] (3.100,0.865) rectangle (7.300,3.780);
\node[align=left,text width=3.940cm,font={\linespread{1}\fontsize{10}{13.50}\selectfont}] at (5.200,2.323) {（1）只要客观上达到“数额较大”，就定罪};
\draw[line width=.4pt] (7.300,0.865) rectangle (16.000,3.780);
\node[align=left,text width=8.440cm,font={\linespread{1}\fontsize{10}{13.50}\selectfont}] at (11.650,2.323) {甲向局长乙送一盒茶叶（藏有3万元现金），乙以为只是一盒普通茶叶便收下。有意见认为，乙构成受贿罪。然而，乙主观上没有收受3万元现金的故意，不构成受贿罪};
\draw[line width=.4pt] (3.100,3.780) rectangle (7.300,7.995);
\node[align=left,text width=3.940cm,font={\linespread{1}\fontsize{10}{13.50}\selectfont}] at (5.200,5.888) {（2）只要客观上达到“数额（特别）巨大”，就适用“数额（特别）巨大”的法定刑};
\draw[line width=.4pt] (7.300,3.780) rectangle (16.000,7.995);
\node[align=left,text width=8.440cm,font={\linespread{1}\fontsize{10}{13.50}\selectfont}] at (11.650,5.888) {甲向局长乙送了一张银行卡，乙问：“卡中有多少钱？”甲担心乙拒收，便答：“3万元”（实际有20万元）。乙收下。有意见认为，对乙应适用“数额巨大”的法定刑。然而，乙主观上只有收受3万元的故意，没有收受20万元的故意，只能适用“数额较大”的法定刑，不能适用“数额巨大”的法定刑};
\draw[line width=.4pt] (0.000,7.995) rectangle (3.100,17.325);
\node[align=left,text width=2.840cm,font={\linespread{1}\fontsize{10}{13.50}\selectfont}] at (1.550,12.660) {误解二：只要客观上不具备数额，就不定罪量刑};
\draw[line width=.4pt] (3.100,7.995) rectangle (7.300,13.610);
\node[align=left,text width=3.940cm,font={\linespread{1}\fontsize{10}{13.50}\selectfont}] at (5.200,10.803) {（1）只要客观上未达到“数额较大”，就不定罪};
\draw[line width=.4pt] (7.300,7.995) rectangle (16.000,13.610);
\node[align=left,text width=8.440cm,font={\linespread{1}\fontsize{10}{13.50}\selectfont}] at (11.650,10.803) {甲真心答应向局长乙送3万元现金，乙办成事后，甲又只想送一千元，乙发现只有1千元。有意见认为，乙只收到1千元，不构成受贿罪。然而，甲开始有送3万元的意思，乙有收受3万元的故意，乙便构成受贿罪，没有收到3万元，构成受贿罪未遂。甲构成行贿罪中止。\\{}注意：如果甲一开始就不想送3万元，欺骗乙“会送3万元”，只送了1千元，那么，甲不构成行贿罪，乙不构成受贿罪，因为乙没有收到数额较大财物的可能性};
\draw[line width=.4pt] (3.100,13.610) rectangle (7.300,17.325);
\node[align=left,text width=3.940cm,font={\linespread{1}\fontsize{10}{13.50}\selectfont}] at (5.200,15.467) {（2）只要客观上未达到“数额（特别）巨大”，就不适用“数额（特别）巨大”的法定刑};
\draw[line width=.4pt] (7.300,13.610) rectangle (16.000,17.325);
\node[align=left,text width=8.440cm,font={\linespread{1}\fontsize{10}{13.50}\selectfont}] at (11.650,15.467) {甲真心答应向局长乙送100万元现金，后甲又后悔，只送了3万元。有意见认为，对乙只适用“数额较大”的法定刑。然而，乙应构成收受“数额巨大财物”的未遂，同时构成收受“数额较大财物”的既遂，想象竞合，择一重刑论处};
\end{tikzpicture}
